\section{Introduction}\label{sec:introduction}\label{altord:section}

Let $f:X\to Z$ be a surjective morphism with connected fibres between
smooth projective varieties. The pluricanonical forms on $X$ measure
its birational complexity. Iitaka's subadditivity problem asks whether
the forms on the base and on the geometric generic fibre always
contribute additively to that complexity. More precisely, the Kodaira
dimension $\kappa(X)$ is the maximum dimension of the images of the
complete linear systems $|mK_X|$, for positive integers $m$; it is
$-\infty$ if every such system is empty. If $F$ is the geometric
generic fibre of $f$, the conjecture $C_{n,m}$ asks for
\[
                 \kappa(X)\geq\kappa(F)+\kappa(Z),
                 \qquad n=\dim X,\quad m=\dim Z.
\]
We use $(-\infty)+a=-\infty$ and $\kappa(\mathrm{point})=0$.

The problem belongs to Iitaka's theory of dimensions of divisors
\cite{IitakaOriginal}. Its development has repeatedly connected the
growth of pluricanonical systems with positivity on the base.
Viehweg's weak-positivity method gives addition when the base is of
general type \cite{ViehwegAdditivity}. Kawamata proved addition over
curves and when the general fibre has a good minimal model
\cite{KawamataCurves,KawamataKodaira}; Koll\'ar established the case
of general-type fibres \cite{Kollar87}. The log-general-type-fibre
theorem of Kov\'acs--Patakfalvi and its formulation by Hashizume
provide the addition theorem used below
\cite{KovacsPatakfalvi,Hashizume}.

Birkar proved the conjecture in total dimension at most six
\cite{BirkarIitakaSix}. Using canonical-bundle and nonvanishing
methods, Chang extended this to total dimension at most seven under
the additional assumption $\kappa(X)\geq0$
\cite[Theorem~1.8]{ChangNonvanishing}.
Other advances include Cao--P\u{a}un's theorem for projective klt pairs over abelian varieties,
Hacon--Popa--Schnell's result for bases of maximal Albanese dimension,
and Cao's theorem for projective klt pairs over surfaces
\cite{CaoPaunAbelian,HaconPopaSchnellAbelian,CaoSurfaces}.
These results show how information about the base and positivity of
direct images can control the total pluricanonical system.
Tsuji's preprint states ordinary
subadditivity in Theorem~1.13, while Maehara's preprint states a
broader variation inequality in Section~7, Theorem~8
\cite{TsujiDirectImages,MaeharaIitaka}. Those preprint claims are not
inputs to our proof. The argument here is a projective proof through
a Hodge line and a canonical bundle formula.

\begin{theorem}[Ordinary Iitaka subadditivity]\label{thm:ordinary}\label{altord:ordinary-subadditivity}
Let $k$ be an algebraically closed field of characteristic zero, and
let $f:X\to Z$ be a surjective projective morphism with connected
fibres between smooth connected projective $k$-varieties. If $F$ is
its geometric generic fibre, then
\[
                     \kappa(X)\geq\kappa(F)+\kappa(Z).
\]
\end{theorem}

Thus the ordinary Iitaka conjecture has a positive answer in this
setting. The main step in the proof is a positivity statement for
the entire highest Hodge line of a rationally polarized variation.
We first describe that statement, then explain how the canonical
bundle formula brings it into the subadditivity problem.

Let $Y$ be a smooth connected projective complex variety and $U$ a
dense open subset. A rationally polarized integral pure variation
$\mathbb V$ on $U$ consists of a rational local system with a
monodromy-invariant lattice, a Hodge filtration of fixed weight, and
a rational flat polarization. Suppose its highest nonzero
filtration step $F^p\mathbb V_{\mathcal O}$ is a line, so
$F^{p+1}=0$ and $\operatorname{rank}F^p=1$. On a simply connected
open set, a flat marking identifies this line with a holomorphic
map to the projective space of a fixed reference fibre. We call
this the \emph{highest-line map}. The \emph{full period map}
records every step of the Hodge filtration.

The highest line has a rational parabolic extension across a smooth
normal-crossing compactification. This is the extension whose pullback,
after a finite level cover and resolution making local monodromy
unipotent, is the Deligne--Schmid highest line. The rational weights
retain the finite parts of the original local monodromies.

\begin{theorem}[Projective adjoint positivity]\label{thm:adjoint}\label{altord:p:observation}
Let $Y$ be a smooth connected projective complex variety, let
$U\subset Y$ be a dense Zariski open subset, and let $\mathbb V$
be a rationally polarized integral pure variation of Hodge structure
on $U$ whose entire highest nonzero Hodge filtration step is a line.
Let $M\in\operatorname{Pic}(Y)\otimes\Q$. Assume that, after a smooth
projective birational modification of $Y$ if necessary, there are a
smooth projective alteration $\tau:\widehat Y\to Y$ and a rational
number $c>0$ such that the pulled-back variation has unipotent
monodromy along a simple-normal-crossing boundary and
\[
                     \tau^*M\sim_{\Q}cJ_{\widehat Y},
\]
where $J_{\widehat Y}$ is its extended highest line.
After further smooth projective birational modification, there are a
surjective morphism $p:Y\to S$ with connected fibres onto a smooth
projective variety and a nef rational line $L$ on $S$ such that
\[
 M\sim_{\Q}p^*L,\qquad
 K_S+jL\text{ is big for all sufficiently large integers }j.
\]
If $S$ is a point, the conclusion is $M\sim_{\Q}0$.
\end{theorem}

The two period maps have separate jobs. The full period image provides
the projective base on which an adjoint divisor can be big. The
numerical dimension of $L$ is measured by the highest-line map,
whose rank can be smaller. The argument must therefore remove the
boundary using the rank of the line, without identifying it with
the dimension of the full period image.

\subsection{The boundary obstruction and its resolution}

The algebraicity theorem of Bakker--Brunebarbe--Tsimerman constructs
the full period image as an algebraic variety
\cite[Theorem~1.1]{BakkerBrunebarbeTsimerman}. At a neat level,
Brunebarbe--Cadorel's theorem gives logarithmic general type when
the full period map is generically immersive
\cite[Theorem~1.1]{BrunebarbeCadorel}. A finite quotient returns to
the original level. This introduces a second divisor to be removed:
besides the boundary with nonzero monodromy logarithm, there is
divisorial ramification of the finite quotient. Ordinary adjoint
positivity requires control of both.

We first retain the smallest rational subvariation containing the
highest line. This is the pure transcendental-part construction
of Bakker--Filipazzi--Mauri--Tsimerman
\cite[Definition~2.9]{BFMT}. Rational minimality has a useful
consequence: a connected rational normal subgroup of monodromy
that fixes the line pointwise must be trivial. Together with
Andr\'e's monodromy normality theorem \cite[\S5]{AndreMonodromy}
and the fixed-part results of Griffiths, Deligne, and Schmid
\cite[added note to~4.2.6 and~4.2.8(i),(ii)]{DeligneHodgeII}
\cite[Theorem~7.22 and Corollary~7.23]{Schmid},
it also detects a flat automorphism
that is scalar on all highest lines and preserves one full period.
Such an automorphism preserves every full period.

The local analytic step concerns a family of holomorphic maps
approaching a divisor. If its limiting tangential rank equals
the full rank in the interior, then its nearby image is already
contained in the limiting image. At a nilpotent degeneration,
the limiting line is in the kernel of the monodromy logarithm;
rational minimality rules out equality of ranks. At a divisor
fixed by finite inertia, the limiting lines lie in one eigenspace.
The same rank principle and the fixed-part theorem would make
the inertia element fix the entire full period image.

For this last contradiction we use the normalization of the
\emph{actual} period image. The finite group acts effectively on
its function field. Equality of the full periods therefore forces
the inertia element to be the identity. A further finite cover
of the image could have a nontrivial automorphism acting trivially
on every period, so connected-fibre factorization is taken only
after this argument. Sections~\ref{sec:rational}--\ref{sec:rankloss}
give the rational, analytic, and finite-quotient parts of this rank
loss in that order.

Once the ranks drop, elementary section counting removes the
boundary and divisorial ramification. If the highest-line map has
generic rank $r$, the ample
perturbation in the interior supplies a polynomial of degree $r$
in a large multiple of $L$, whereas the restriction polynomials
on those divisors have degree at most $r-1$. Effective divisors
exceptional over the quotient do not change its section growth.
This yields the adjoint positivity in Theorem~\ref{thm:adjoint}.

\subsection{From the period theorem to ordinary subadditivity}

For the nonvacuous case of Theorem~\ref{thm:ordinary}, put
$d=\kappa(F)\geq0$ and assume $\kappa(Z)\geq0$. The relative Iitaka
fibration factors a birational model of $X$ through a variety $Y$
with $\dim(Y/Z)=d$ and Kodaira-dimension-zero generic fibre over
$Y$. The canonical bundle formula separates the singular-fibre
contribution from a nef moduli divisor:
\[
 K_{\widetilde X}\sim_{\Q}
 g^*(K_Y+B+M)+R,\qquad \widetilde X\xrightarrow{g}Y\xrightarrow{q}Z.
\]
Here $B$ is an effective klt boundary, and the residual divisor $R$
has the direct-image and exceptionality properties needed to compare
complete pluricanonical systems. This is the strategy of
Fujino--Mori's canonical bundle formula, with the birational
discriminant and moduli theory of Ambro
\cite[Theorem~4.5]{FujinoMori}
\cite{AmbroBoundary,AmbroModuli}.

The least-index canonical root cover of the Kodaira-zero fibre
has geometric genus one. Its unique top form is therefore the
entire highest line of a rational variation, even if the character
of the cyclic covering group is not rational. The Hodge realization
of $M$ matches its extension on an alteration, not only its
restriction to the open base. These facts place $M$ within
Theorem~\ref{thm:adjoint}. The auxiliary sub-boundary in this formula
can be negative on the generic fibre; the argument uses nefness
and the Hodge realization, not a semiampleness assertion requiring
an effective generic boundary.

The section comparison gives
\[
 \kappa(X)\geq\kappa(Y,K_Y+B+M),\qquad
 (K_Y+B+M)|_{Y_{\bar\eta}}\text{ big},
\]
where $Y_{\bar\eta}$ is the geometric generic fibre of $q$.
The period projection writes $M=p^*L$ for a morphism $p:Y\to S$.
Write $\mu:\widetilde X\to X$ for the birational modification
of the original source in the relative Iitaka construction.
Figure~\ref{fig:two-bases} separates the two maps from $Y$: their
bases need not map to one another.

\begin{figure}[ht]
\centering
\begin{tikzcd}[column sep=large,row sep=large]
\widetilde X \arrow[r,"g"] \arrow[dr,"f\mu"'] &
Y \arrow[d,"q"] \arrow[r,"p"] & S \\
& Z &
\end{tikzcd}
\caption{The relative Iitaka map $g$ has Kodaira-dimension-zero
generic fibre. The original base map $q$ and the independent period
map $p$ start from the same variety $Y$. The divisor $M=p^*L$
comes from $S$, while the desired addition inequality is measured
over $Z$.}
\label{fig:two-bases}
\end{figure}


A product-map argument first gives log pluricanonical nonvanishing on a general fibre
of $p$. Fujino's twisted weak positivity then turns the big
divisor $K_S+jL-H$, for an ample $H$ and large $j$, into an
effective divisor linearly equivalent to
$K_Y+B+jM-p^*H$ \cite[Theorem~1.1]{FujinoWeakPositivity}.
Interpolating this divisor with an effective klt log canonical
class reduces the desired inequality to the established addition
theorem for log-general-type fibres
\cite[Theorem~2.11]{Hashizume}. Scalar extension of each
pluricanonical space transfers the result from $\C$ to any
algebraically closed field of characteristic zero.

There are two further proofs of the Hodge-line boundary estimate
in the projective setting. An integral tensor replacement makes
monodromy faithful on the moving highest line; projection to a
fixed quotient by a monodromy logarithm then produces transverse
loops on which the line is constant. A second calculation obtains
the restriction-intersection bound directly from Hodge curvature:
smooth Thom forms concentrate on a boundary divisor, and a product
Poincar\'e estimate controls their limits at crossings. These
two calculations share the integral tensor replacement and the
transverse-loop criterion; the second supplies a different proof
of the boundary restriction-intersection step. These
arguments retain the same whole-top-line hypothesis but isolate
different mechanisms behind the boundary estimate. In the main
route, rational minimality and the tangential rank principle treat
the finite quotient on the actual period image directly.

\subsection{Organization}

Section~\ref{sec:rational} constructs the rationally minimal variation and proves the
fixed-period criterion. Section~\ref{sec:limits} compares numerical rank with
boundary limits and proves the tangential rank principle.
Section~\ref{sec:rankloss} applies these results to boundary monodromy and divisorial
ramification. Section~\ref{sec:quotient} constructs the actual period quotient and
completes the proof of adjoint positivity. Section~\ref{sec:canonical} gives the
canonical bundle data, the least-index root cover, and the
section-growth comparison. Section~\ref{sec:ordinary} proves ordinary subadditivity
and gives the scalar-extension argument in characteristic zero.
Section~\ref{p1logalt:section} proves the integral tensor replacement and the alternative
adjoint argument by transverse loops. Section~\ref{p1varhd:section} proves the
restriction-intersection estimate using Thom forms, including the
crossing calculation.

\subsection{Conventions}
An algebraic fibre space is a surjective projective morphism with
connected fibres. All varieties in the Hodge-theoretic argument are
over $\C$. Divisors and line-bundle identities are rational unless
otherwise specified, and section comparisons are taken in degrees
clearing all denominators. We write $A\sim_{\Q}B$ when a positive
integral multiple of $A-B$ is principal. An effective rational class
means a class rationally linearly equivalent to an effective divisor.
The positive and negative parts of a rational divisor have disjoint
support. Numerical dimension is used only for nef divisors and is
defined in Section~\ref{sec:limits}.
